% ==== AUTO-GENERADA — no editar a mano (regenerar con generate_thesis_tables.py) ====
\begin{table}[htbp]
\centering
\caption[Predicción a $N$ grande (heavy-hex)]{Predicción entre tamaños (sobre casos no observados) en heavy-hex: la UnifiedMPNN entrenada con $N$ pequeños predice ángulos para $N$ grandes.}
\label{tab:auto_heavy_hex_large_n}
\begin{tabular}{rrrr}
\toprule
$N$ & $|\Delta E|$ & gap & $\Delta E/\mathrm{gap}$ (\%) \\
\midrule
21 & 0,5194 & 6,1950 & 8,38 \\
22 & 0,0093 & 6,1859 & 0,15 \\
24 & 0,0160 & 5,6703 & 0,29 \\
26 & 0,0396 & 3,6580 & 2,15 \\
30 & 0,0204 & 5,6357 & 0,37 \\
32 & 1,0073 & 3,6279 & 88,70 \\
40 & 0,0215 & 5,6013 & 0,40 \\
50 & 0,9189 & 3,5811 & 28,56 \\
60 & 1,9493 & 3,5675 & 54,67 \\
\bottomrule
\end{tabular}
\\[2pt]
\footnotesize Evaluado sobre la grilla de $h$ canónica del régimen de extrapolación, $h \in \{2,5; 3,0; 3,5; 4,0; 4,5; 5,0\}$, restringida al régimen válido ($h \geq h_{\min}$) de cada $N$; por eso los $N$ mayores contribuyen con menos puntos. $|\Delta E|$ y gap son medias sobre esos puntos, en unidades de $J = 1$. La columna $\Delta E/\mathrm{gap}$ es la media de los cocientes \emph{por punto} (no el cociente de las dos columnas anteriores): cerca de la frontera algún punto tiene el gap estrecho y su cociente individual es grande, de modo que la media de cocientes puede superar al cociente de las medias. La métrica primaria es $|\Delta E|$; $\Delta E/\mathrm{gap}$ es el criterio de clasificación normalizado. Fuente: elaboración propia.
\end{table}
